% arara: lualatex_git
\documentclass[4x6,food]{recipe-card}

\begin{document}

\begin{recipe}{Weeknight Pomodoro (Demo)}
\setRecipeType{Primi}
\begin{ingredients}
\group{Sauce}
\ingredient{2 tbsp}{Olive Oil}
\ingredient{2 cloves}{Garlic, Sliced}
\ingredient{400 g}{Crushed Tomatoes}
\ingredient{1 pinch}{Chili Flakes}
\group{Pasta}
\ingredient{280 g}{Spaghetti}
\ingredient{to taste}{Salt and Black Pepper}
\group{Finish}
\ingredient{as needed}{Fresh Basil}
\ingredient{to serve}{Grated Parmigiano}
\end{ingredients}
\begin{process}
\step{Warm olive oil in a pan, then gently cook garlic and chili for 1 minute.}
\step{Add tomatoes, season, and simmer until slightly thickened.}
\step{Cook spaghetti in salted water until al dente, then reserve some pasta water.}
\step{Toss pasta with sauce, adding a splash of pasta water as needed.}
\step{Finish with basil and parmigiano before serving.}
\end{process}

\vfill
\begin{recipeinfo}
Quick pantry-friendly first course to test food-mode layout and category bars.
\end{recipeinfo}

\back
\drawBackMirrorRecipeBar
\begin{center}
    \footnotesize\textsc{\color{recipeRed} Notes}
\end{center}
\scriptsize
Designed as a lightweight non-cocktail fixture for testing:
\begin{itemize}
\setlength{\itemsep}{1pt}
\setlength{\parskip}{0pt}
\setlength{\topsep}{2pt}
\item Uses class option \texttt{[4x6,food]}.
\item Uses \texttt{\textbackslash setRecipeType\{Primi\}} instead of \texttt{\textbackslash setBaseSpirit}.
\item Omits cocktail glass/method icon row on the front side.
\end{itemize}

\end{recipe}

\end{document}
